\section{Results}
\label{sec:results}

We present results in the same order as the paper's claims. We start with null
diagnostics that speak directly to Stage~A screening, then move to synthetic
ranking quality, runtime, real-data behavior, and ablations.

\subsection{Null diagnostics}
\label{sec:experiments-null}

\paragraph{Selection bias under a complete global null.}

\begin{figure}[t]
  \centering
  \includegraphics[width=0.95\linewidth]{figures/selection_bias_demo.png}
  \caption{Under a complete global null, greedy CART-style split search favors
  high-cardinality features, whereas Stage~A screening returns a stump after
  no feature clears the Bonferroni-adjusted threshold.}
  \label{fig:selection-bias-demo}
\end{figure}

\Cref{fig:selection-bias-demo} illustrates the split-selection bias that
motivates the paper. Under $X \perp\!\!\!\perp Y$, greedy CART-style split
optimization selects the continuous feature most often because it has the
largest set of admissible thresholds. Stage~A screening removes that effect:
with the same null data, no feature clears the Bonferroni-adjusted threshold
and the fitted tree remains a stump.

\paragraph{Calibration of Stage~A p-values.}

\begin{figure}[t]
  \centering
  \includegraphics[width=0.95\linewidth]{figures/calibration_pvalue_ecdf.png}
  \caption{Empirical CDF and histogram of Stage~A $+1$ permutation p-values
  under a complete global null ($n=200$, $B=199$, 10{,}000 simulations).
  Fixed-$B$ curves remain at or below the Uniform$(0,1)$ diagonal. The
  adaptive-stopping curve is shown only as a descriptive comparison.}
  \label{fig:pvalue-calibration}
\end{figure}

\begin{table}[t]
  \centering
  \caption{Empirical rejection rate under a complete global null at
  $\alpha=0.05$, across permutation budgets $B$ and stopping rules.}
  \label{tab:calibration}
  \small
  \begin{tabular}{llcc}
    \toprule
    Selector & Stopping & $B$ & Rejection rate \\
    \midrule
    MC & fixed-$B$ & 49  & 0.041 \\
    MC & fixed-$B$ & 99  & 0.046 \\
    MC & fixed-$B$ & 199 & 0.055 \\
    MC & fixed-$B$ & 499 & 0.050 \\
    MC & fixed-$B$ & 999 & 0.050 \\
    PC & fixed-$B$ & 199 & 0.049 \\
    MC & adaptive  & 199 & 0.047 \\
    \bottomrule
  \end{tabular}
\end{table}

Under a complete global null with $n=200$ observations, the fixed-$B$ $+1$
permutation p-values are super-uniform, matching the theoretical calibration
target in \cref{sec:theory}. \Cref{tab:calibration} shows empirical rejection
rates close to the nominal 0.05 level across permutation budgets. The
adaptive-stopping variant tracks the fixed-$B$ behavior closely in these runs,
but we report it here only descriptively rather than as part of the
theorem-backed calibration story.

\paragraph{Power and resolution.}

\begin{figure}[t]
  \centering
  \includegraphics[width=0.48\linewidth]{figures/ablation_power.png}
  \caption{Power curve for Stage~A screening at $\alpha=0.05$. Both stopping
  rules control null rejection and achieve similar power once the signal is
  nonzero.}
  \label{fig:power}
\end{figure}

\Cref{fig:power} shows that fixed-$B$ and adaptive stopping have nearly the
same power once signal is present. The main practical lesson is not a
difference in power curves but a difference in resolution: if $B$ is too
small, the minimum attainable p-value $1/(B+1)$ can exceed the target level.
For example, with $B=49$ the test cannot reject at $\alpha=0.01$ regardless of
signal strength. This is a fixed-$B$ design constraint, not a pathology of the
selector statistic.

\subsection{Synthetic ground-truth ranking quality}
\label{sec:experiments-synthetic}

\begin{figure}[t]
  \centering
  \includegraphics[width=\linewidth]{figures/synthetic_cd_precision_at_10_clean.png}
  \caption{Critical-difference diagram for precision@10 on synthetic datasets
  with known informative features. Methods connected by a gray bar are not
  significantly different at $\alpha=0.05$.}
  \label{fig:synthetic-ranking}
\end{figure}

On synthetic datasets with known informative feature sets,
\cref{fig:synthetic-ranking} places conditional-inference forests in the main
competitive group rather than at either extreme. The method performs well on
structured regimes with correlated or nonlinear signal, where deeper trees can
recover interactions and redundancy. Its weak point is the weak-signal regime:
when marginal evidence is diffuse, conservative Stage~A screening prunes
features that are real but individually modest, which lowers precision@10
relative to more permissive ensembles.

This synthetic split is important for interpreting the rest of the paper. The
method is not uniformly weak or uniformly strong. It is strongest when signal
is structured enough for the forest to benefit from recursive partitioning, and
least competitive when success depends on admitting many weak candidates early.

\subsection{Runtime and scalability}
\label{sec:experiments-runtime}

\begin{figure}[t]
  \centering
  \includegraphics[width=0.75\linewidth]{figures/ablation_scaling.png}
  \caption{Runtime scaling. Left: runtime grows roughly linearly in $n$ at
  fixed $p$. Right: runtime is comparatively flat over the tested range of $p$,
  but this should not be read as dimension-independence of the algorithm.}
  \label{fig:runtime}
\end{figure}

Permutation testing adds a clear computational premium, but the runtime
behavior is still interpretable and manageable for offline ranking.
\Cref{fig:runtime} shows roughly linear growth in $n$ over the tested range.
The observed weak dependence on $p$ is more subtle: it reflects the realized
work on these benchmarks rather than ambient dimension disappearing from the
algorithm. The implementation still pays feature- and threshold-dependent
preprocessing costs, but the dominant permutation loops are capped effectively
enough that those costs do not control the wall-clock profile in this slice of
the benchmark.

Early stopping matters materially for usability. In the easy classification
setting used in the ablation runs, disabling early stopping increases runtime
by roughly a factor of four with little change in ranking quality. That makes
early stopping one of the few engineering choices that is clearly worth keeping
even when the paper stays conservative about its inferential interpretation.

\subsection{Real-data benchmark}
\label{sec:experiments-real}

\begin{figure}[t]
  \centering
  \includegraphics[width=\linewidth]{figures/clf_cd_balanced_accuracy_clean.png}
  \caption{Exploratory critical-difference diagram for classification,
  averaging across downstream models and values of $k$. This is a descriptive
  robustness summary rather than the pre-specified primary-endpoint analysis.}
  \label{fig:real-clf}
\end{figure}

\begin{figure}[t]
  \centering
  \includegraphics[width=\linewidth]{figures/reg_cd_r2_clean.png}
  \caption{Exploratory critical-difference diagram for regression, again
  averaging across downstream models and values of $k$.}
  \label{fig:real-reg}
\end{figure}

The real-data figures in \cref{fig:real-clf,fig:real-reg} are intentionally
presented as descriptive summaries. They average across secondary choices of
downstream model and $k$, so their role is to show whether the qualitative
story is stable rather than to carry confirmatory pairwise claims.

That qualitative story is consistent with the synthetic benchmark. The method
is most competitive on datasets where correlation structure or nonlinear
interactions make recursive screening useful, and less competitive in sparse,
weak-signal, or very high-dimensional regimes where conservative Stage~A
screening suppresses depth before the forest can accumulate weaker discoveries.
The finalized primary-endpoint tables should therefore be read as the source of
exact real-data claims, with these figures serving as context for the broader
pattern.

\subsection{Ablation study}
\label{sec:experiments-ablation}

The ablations are easiest to read around one question: which knobs change
ranking quality because they change the \emph{statistical strictness} of the
screening rule, and which mainly save compute?

\paragraph{Screening strictness.}

\begin{figure}[t]
  \centering
  \includegraphics[width=0.85\linewidth]{figures/ablation_alpha_depth.png}
  \caption{Precision@10 versus $\alpha_{\mathrm{sel}}$. Strong-signal datasets
  change little across the range; weak-signal datasets benefit from more
  permissive screening.}
  \label{fig:ablation-alpha}
\end{figure}

\Cref{fig:ablation-alpha} shows that strong-signal datasets are largely
insensitive to $\alpha_{\mathrm{sel}}$, whereas weak-signal datasets benefit
from more permissive screening. The same pattern appears even more strongly in
the Bonferroni toggle: in the current weak-signal runs, removing the
correction makes trees noticeably deeper and improves precision@10. The exact
magnitudes may move as the benchmark fills in, but the structural tradeoff is
already clear. Conservative nodewise control improves screening strictness,
yet it also limits the forest mechanism that discovers weaker features deeper
in the tree.

\paragraph{Forest size.}

\begin{figure}[t]
  \centering
  \includegraphics[width=0.85\linewidth]{figures/ablation_n_estimators.png}
  \caption{Precision@10 versus number of estimators. Performance typically
  saturates by roughly 25 trees, with weaker-signal regimes showing the most
  variability.}
  \label{fig:ablation-nest}
\end{figure}

Increasing the number of trees helps quickly and then saturates
(\cref{fig:ablation-nest}). On structured synthetic datasets, most of the gain
arrives by about 25 trees. Weak-signal settings are less monotone: adding more
trees can average away sparse discoveries when individual trees disagree on
which modest signals are worth keeping.

\paragraph{Permutation budget.}

\begin{figure}[t]
  \centering
  \includegraphics[width=0.85\linewidth]{figures/ablation_nresamples.png}
  \caption{Precision@10 versus permutation budget. Ranking quality changes
  little once the budget is large enough to resolve the relevant thresholds.}
  \label{fig:ablation-nresamples}
\end{figure}

\Cref{fig:ablation-nresamples} shows little ranking benefit beyond modest
permutation budgets once early stopping is enabled. This does not mean that
resampling is irrelevant; rather, it means that after the budget is large
enough to resolve the practically relevant decisions, additional permutations
mainly increase cost.

\paragraph{Bootstrap and feature subsampling.}

\begin{figure}[t]
  \centering
  \includegraphics[width=\linewidth]{figures/ablation_bootstrap.png}
  \caption{Precision@10 under different bootstrap and feature-subsampling
  strategies. The best choice depends on whether the dataset is dominated by
  easy signal or by correlated confounding structure.}
  \label{fig:ablation-bootstrap}
\end{figure}

Bootstrap and feature subsampling interact strongly with dataset structure
(\cref{fig:ablation-bootstrap}). On easy strong-signal data, no-bootstrap
training can be best because bagging only injects ranking variance. On
confounded or correlated data, feature subsampling helps by forcing different
trees to consider different parts of the feature space instead of repeatedly
locking onto the same dominant variables.

\paragraph{Noise robustness.}
When 1000 pure-noise features are injected, the method still recovers many
informative variables on strong-signal data, but it degrades noticeably on
moderate or weak signal. The important point is not that noise breaks the
method completely; it is that high-dimensional noise amplifies the same
strictness--depth tradeoff seen elsewhere in the paper. Once the Stage~A gate
becomes too conservative, extra trees in the forest cannot compensate for the
lack of depth.
